\documentclass[conference]{IEEEtran}
\usepackage{cite}
\usepackage{amsmath,amssymb,amsthm}
\usepackage[hidelinks]{hyperref}

\newtheorem{theorem}{Theorem}
\newtheorem{lemma}[theorem]{Lemma}
\newtheorem{corollary}[theorem]{Corollary}
\newcommand{\Z}{\mathbb{Z}}
\newcommand{\one}{\mathbf{1}}
\newcommand{\qminus}{q_-}
\newcommand{\qplus}{q_+}

\begin{document}

\title{Asymptotic Growth of the Number of Rubik's Snake Shapes}
\author{\IEEEauthorblockN{Dmytro Fedoriaka}
\IEEEauthorblockA{\textit{University of Washington}\\
Seattle, Washington, USA\\
fedimser@cs.washington.edu}}
\maketitle

\begin{abstract}
Let $S_n$ count the collision-free configurations of an oriented
$n$-wedge Rubik's Snake with quarter-turn joints. We prove that
$\mu=\lim_{n\to\infty}S_n^{1/(n-1)}$ exists and give computer-assisted
bounds
$3.400034903<\mu\leq3.667542939$.
The lower bound comes from uniquely decomposable paths confined to
successive disjoint slabs; the upper bound comes from a finite graph
enforcing all collision constraints within a prescribed window.
Finite enumerations are converted into exact integer certificates,
rather than estimates of eigenvalues or polynomial roots. The same
constructions give the pointwise bounds
$0.014(3.4000)^n<S_n$ for $n\geq4$ and
$S_n<2.95\cdot10^{11}(3.6855)^n$ for $n\geq10$.
Reproducible notebooks implement each bound with adjustable truncation
parameters. No constant-prefactor asymptotic or ratio limit is assumed.
\end{abstract}

\section{Introduction}

Rubik's Snake is a chain of congruent right isosceles triangular
prisms, called \emph{wedges}. Two square faces of each internal wedge
are joined to its neighbors. Rotating these joints changes the shape
of the chain. We consider only multiples of a quarter turn and forbid
intersections of positive volume. Boundary contact is allowed.
The familiar toy has 24 wedges, but this counting problem is defined
for every positive length.

We distinguish the two ends of the chain and count configurations up
to translation and rotation, not up to reflection or reversal.
Equivalently, after fixing the first wedge, we count valid words of
joint rotations. The resulting sequence $S_n$ is recorded through
$n=28$ in OEIS A375865~\cite{oeisA375865}. Its value
$S_{24}=13\,446\,591\,920\,995$ agrees with Aylett's corrected
enumeration~\cite{aylett2011}.

Rotation-sequence descriptions and geometric design methods have
been developed by Li, Hou, and Bishop~\cite{li2020}; related work
addresses palindromic and periodic snakes~\cite{palindromic2021},
geometric structures~\cite{chen2023}, and planar
configurations~\cite{grotto2024}. Our subject is instead the
exponential growth of the unrestricted three-dimensional count.

The model resembles a self-avoiding lattice walk, but differs in
two essential ways: the direction must turn through a right angle
at every cube, and two complementary wedges may occupy the same
cube. We retain both features throughout. Submultiplicativity,
as in the theory of connective constants~\cite{hammersley1957},
proves existence of a single exponential growth constant.
For upper bounds we adapt finite-memory and overlapping-path
methods~\cite{alm1993,noonan1998,ponitz2000,he2025}.
For lower bounds we construct a family of concatenable slab paths
and remove the ambiguity in their decomposition.

The structural arguments below are independent of computation.
The reported numerical bounds require only finite, exhaustive
enumerations followed by exact inequalities. Section~\ref{sec:reproduce}
identifies a separate notebook for each computation.

\section{Definitions and geometry}\label{sec:geometry}

\subsection{Wedges and rotation words}

Write $e_x,e_y,e_z$ for the coordinate unit vectors and
$\mathcal D=\{\pm e_x,\pm e_y,\pm e_z\}$.
The cube centered at $p\in\Z^3$ is
$Q_p=p+[-1/2,1/2]^3$.
For perpendicular $a,b\in\mathcal D$, define the wedge
\begin{equation}\label{eq:wedge}
 W(p;a,b)=\{u\in Q_p:(a+b)\mathbin{\cdot}(u-p)\geq0\}.
\end{equation}
It has volume $1/2$. Its two square faces are the cube faces
with outward normals $a$ and $b$; its remaining faces are two
triangles and a rectangle. The order of $a,b$ specifies which
square face is the entrance and which is the exit.

A \emph{formula} for $n$ wedges is a word
$r_1\cdots r_{n-1}$ over $\{0,1,2,3\}$.
The symbol records a relative rotation by that many quarter turns,
with zero denoting the untwisted, zigzag configuration. This is the
usual four-symbol notation~\cite{soul2022,li2020}.
For completeness the following coordinate convention specifies it
without a physical toy. Put $p_0=0$, $d_0=e_y$, $d_1=e_x$, and
\begin{equation}\label{eq:directions}
 d_{i+1}=R_{d_i}(r_i\pi/2)d_{i-1},\qquad
 p_i=p_{i-1}+d_i\quad(1\leq i<n),
\end{equation}
where $R_d(\theta)$ is the right-handed rotation about the oriented
axis $d$. Wedge $i$, with indices starting at zero, is
$W(p_i;-d_i,d_{i+1})$. The vectors $d_0$ and $d_n$ describe
the free faces at the two ends, not extra wedges.

Thus the first displacement is fixed, consecutive displacements
are perpendicular, and each of the four possible next directions
corresponds to exactly one rotation symbol. The final symbol still
matters: it determines the last wedge's orientation even though
it does not change any cube center. A center-line walk alone is
therefore not a complete formula.

A formula is \emph{valid} if the interiors of all its wedges are
pairwise disjoint. We define $S_n$ to be the number of valid formulas.
In particular $S_1=1$ and $S_n\leq4^{n-1}$.
The chain order and its distinguished head are retained; configurations
are not identified merely because their occupied unions coincide.
We count static configurations, without imposing a collision-free
continuous motion from a particular starting shape.
Choosing only positive coordinate directions gives two choices
at each joint and never repeats a cube. Hence also
$S_n\geq2^{n-1}$.

\subsection{The exact collision rule}

Wedges in distinct lattice cubes cannot overlap in positive volume.
Within one cube, the complement of $W(p;a,b)$ is $W(p;-a,-b)$,
up to their common boundary. These are the only two wedges that
can share a cube: two half-volume wedges with disjoint interiors
must fill the cube, and their cutting planes must coincide.
Consequently, a cube is either empty, occupied once, or occupied
twice by complementary wedges; a third visit is impossible.

\begin{figure}[t]
\centering
\setlength{\unitlength}{1pt}
\begin{picture}(230,110)
\put(65,12){\framebox(70,70){}}
\put(65,82){\line(1,-1){70}}
\put(111,61){$W$}
\put(76,28){$W^c$}
\put(135,47){\vector(1,0){30}}
\put(171,44){$a$}
\put(100,82){\vector(0,1){20}}
\put(108,99){$b$}
\end{picture}
\caption{A cross-section of a unit cube in the plane of $a$ and
$b$. The triangles extend along the remaining coordinate to
form complementary wedges. The square faces of $W$ have
outward normals $a$ and $b$.}
\label{fig:wedge}
\end{figure}

For a path arriving with step $d$ and leaving with step $e$, its
occupancy is the unordered pair $\{-d,e\}$. At a second visit,
the pair must be the negative of the first pair. This finite rule
replaces continuous collision detection in all our enumerations.
It also explains why forbidding all repeated vertices would count
a strictly smaller family; see Fig.~\ref{fig:wedge}.

\section{Finite enumeration of \texorpdfstring{$S_n$}{Sn}}

A depth-first search starts with the fixed first wedge and tries
all four symbols at each joint. A branch stops immediately when
the new wedge fails the collision rule. Recording the number of
successful prefixes gives all counts through the chosen length.
The implementation stores the occupancy type of each visited cube
and precomputes the 24 ordered wedge orientations and their
rotation transitions. No floating-point geometry is needed.

The existing notebook \texttt{count-shapes.ipynb} describes this
enumeration; the stored sequence extends through 28 wedges.
Selected terms appear in Table~\ref{tab:counts}. Long unrestricted
enumeration is expensive, so our bounds use shorter local windows
and specially constrained paths rather than additional terms of
$S_n$.

\begin{table}[t]
\caption{Selected exact counts of valid formulas~\cite{oeisA375865}.}
\label{tab:counts}
\centering
\begin{tabular}{r|r@{\qquad}r|r}
$n$ & $S_n$ & $n$ & $S_n$\\ \hline
1 & 1 & 8 & 12\,585\\
2 & 4 & 9 & 46\,471\\
3 & 16 & 10 & 172\,226\\
4 & 64 & 12 & 2\,333\,757\\
5 & 241 & 16 & 422\,677\,188\\
6 & 920 & 20 & 75\,663\,420\,126\\
7 & 3\,384 & 24 & 13\,446\,591\,920\,995\\ \hline
\multicolumn{2}{r}{$n=28$:} &
\multicolumn{2}{l}{$S_{28}=2\,377\,810\,831\,870\,022$}
\end{tabular}
\end{table}

\section{Exponential growth}

\begin{lemma}\label{lem:heredity}
Every contiguous subword of a valid formula is valid.
\end{lemma}
\begin{proof}
The subword specifies a contiguous subchain. Deleting the remaining
wedges cannot create an overlap. A rigid rotation and translation
put its first wedge in the prescribed initial orientation, without
altering the relative joint rotations or internal intersections.
\end{proof}

\begin{theorem}\label{thm:existence}
There is a constant $\mu$ such that
\begin{equation}\label{eq:mu}
 \lim_{n\to\infty}S_n^{1/(n-1)}
 =\mu=\inf_{k\geq1}S_{k+1}^{1/k}.
\end{equation}
In particular $S_n=\mu^{n+o(n)}$ and
$\lim_{n\to\infty}S_n^{1/n}=\mu$.
\end{theorem}
\begin{proof}
Set $a_k=S_{k+1}$, with $a_0=1$. Splitting a valid word after $p$
symbols is an injection into pairs of valid words of lengths
$p$ and $q$, by Lemma~\ref{lem:heredity}. Hence
$a_{p+q}\leq a_pa_q$.
For an elementary subadditivity argument, fix $k\geq1$ and write
$N=tk+r$, $0\leq r<k$. Then
$a_N\leq a_k^t a_r$, so
$\limsup_N N^{-1}\log a_N\leq k^{-1}\log a_k$.
Taking the infimum over $k$ gives the reverse inequality to the
one that holds for every term. This proves convergence and
\eqref{eq:mu}.
\end{proof}

This conclusion permits oscillating consecutive ratios
$S_{n+1}/S_n$. It does not assert either convergence of those
ratios or an asymptotic equivalence $S_n\sim C\mu^n$.

\subsection{Lower bounds from slab renewals}\label{sec:lower}

\subsubsection{Separated blocks}
Fix $+e_x$ as a progress direction. For $w\geq0$, a
\emph{slab block} consists of a path of $\ell$ internal edges
starting at $(0,0,0)$, staying in $0\leq x\leq w$, and ending
in the plane $x=w$. It has a prescribed incoming $+e_x$ edge
at its start and an outgoing $+e_x$ edge at its end.
Consecutive directions, including these two boundary directions,
must be perpendicular. All $\ell+1$ wedges in the slab must
satisfy the exact collision rule. Neither external endpoint
is a wedge of the block.

The outgoing edge enters a fresh slab. Translating successive
blocks so that each starts just beyond the preceding slab
therefore gives a valid path. The two natural sizes of a block
are its progress $d=w+1$ and its length $j=\ell+1$ (internal
edges plus the outgoing edge). Write $B_{d,j}$ for the number
of blocks of these sizes, with the starting point fixed.
In particular $B_{d,j}=0$ for $j<2d$: every one of the at
least $d-1$ internal forward steps must be separated from the
next, and from each boundary edge, by a transverse step.

These numbers are computed by a depth-first search with incoming
direction $+e_x$. Before advancing, the chosen outgoing direction
finalizes the current wedge, whose occupancy is checked and then
restored on backtracking. At every endpoint in $x=w$ with a
transverse incoming direction, the search tries the prescribed
exit $+e_x$ and records the block if the final wedge is legal.
Continuing the search computes all lengths through the cutoff
in one traversal. Transverse coordinates have absolute value
at most that cutoff, so the enumeration is finite.

\subsubsection{Why irreducible blocks are necessary}
A path may admit more than one decomposition into arbitrary
slab blocks. Adding counts from different widths as if they
were independently distinguishable building blocks would
therefore overcount.

An internal $+e_x$ edge from level $h$ to level $h+1$ is a
\emph{separator} if all vertices before that edge have $x\leq h$
and all vertices after it have $x\geq h+1$. Cutting at every
such edge gives an \emph{irreducible} block decomposition.
The separator property is determined by the whole path, so
these cuts are unique. Conversely, a concatenation of irreducible
blocks has a separator at every join. An edge inside a factor
is a separator of the concatenation exactly when it is one
of that factor: the other factors lie strictly to its left
or right. Thus the decomposition is uniquely recoverable.

Let $I_{d,j}$ count irreducible blocks, and introduce the formal
power series
\[
 \begin{aligned}
 B(X,Z)&=1+\sum_{d,j\geq1}B_{d,j}X^dZ^j,\\
 I(X,Z)&=\sum_{d,j\geq1}I_{d,j}X^dZ^j.
 \end{aligned}
\]
Progress and length add at joins, giving
\begin{equation}\label{eq:factorization}
 B=1+I+I^2+\cdots=\frac1{1-I}.
\end{equation}
Only finitely many terms contribute to each coefficient, so this
is a formal identity, with no convergence assumption. It gives
the exact recurrence
\begin{equation}\label{eq:irreducible}
 I_{d,j}=B_{d,j}-
 \sum_{a=1}^{d-1}\sum_{b=1}^{j-1}
 I_{a,b}B_{d-a,j-b}.
\end{equation}

The bivariate bookkeeping is important also when truncating the
enumeration. If tables are supplied for $1\leq d\leq D$ through
internal lengths $K_d$, with
$K_1\geq K_2\geq\cdots\geq K_D$, then every coefficient needed
to compute $I_{d,j}$ for $j\leq K_d+1$ is available.
Unknown tails are not substituted into this recurrence as zeros.
After the known irreducible coefficients have been extracted,
we may omit all unknown coefficients, since they are nonnegative.

\subsubsection{From counts to a lower bound}
Let $i_j$ be the number of retained irreducible blocks of length
$j$, summed over progress, and let $L$ be the largest retained
length. Put
\begin{equation}\label{eq:renewal}
 c_0=1,\quad c_k=0\ (k<0),\qquad
 c_k=\sum_{j=1}^L i_jc_{k-j}\quad(k\geq1).
\end{equation}
Unique factorization makes $c_k$ the exact number of
concatenations of length $k$, not just a formal weighted count.

Prepend the fixed first wedge and its $+e_x$ edge to such a
concatenation. Its final outgoing edge reaches a new cube,
where we add a final wedge with, for example, last rotation
zero. The result has $k+2$ wedges. All block wedges are in
disjoint slabs, and the two added wedges are outside them.
Equation~\eqref{eq:directions} recovers the rotation word
uniquely from this direction sequence. Therefore
\begin{equation}\label{eq:lower-pointwise}
 c_{n-2}\leq S_n\qquad(n\geq2).
\end{equation}
This also accounts for the terminal-wedge orientation, which
must not be dropped when passing between paths and formulas.

\begin{theorem}\label{thm:renewal}
For any nonempty finite family of irreducible blocks, let
$\alpha>0$ be the unique solution of
\begin{equation}\label{eq:root}
 \sum_{j=1}^L i_j\alpha^{-j}=1.
\end{equation}
Then $\mu\geq\alpha$. In particular, if $p,d$ are positive
integers and
\begin{equation}\label{eq:lower-certificate}
 p^L-\sum_{j=1}^L i_jp^{L-j}d^j<0,
\end{equation}
then $\mu>p/d$.
\end{theorem}
\begin{proof}
The left side of \eqref{eq:root} is continuous and strictly
decreasing from infinity to zero. Suppose $\mu<\alpha$.
By Theorem~\ref{thm:existence} and \eqref{eq:lower-pointwise},
the series $C=\sum_{k\geq0}c_k\alpha^{-k}$ would converge:
its terms are eventually bounded by a geometric sequence
of ratio less than one. Summing \eqref{eq:renewal} would then
give $C=1+C\sum_j i_j\alpha^{-j}=1+C$, a contradiction.
Finally, \eqref{eq:lower-certificate} says
$\sum_j i_j(p/d)^{-j}>1$, so $p/d<\alpha$.
\end{proof}

For width zero, the blocks are nonempty transverse runs in
a single $yz$-plane. Their axes alternate between $y$ and $z$,
but complementary revisits are allowed. Such a block has no
internal separator, so it is already irreducible. The first
counts, by internal length $\ell=1,2,\ldots$, are
$4,8,16,24,40,72,136,224,\ldots$.
Using lengths through 28 gives $\mu>3.144384190$.
This realizes a plane-by-plane lower-bound strategy related
to Aylett's geometric construction~\cite{aylett2011}.

For the stronger bound we enumerate progress $d=1,2,3$ through
internal lengths $28,20,18$, respectively, and apply
\eqref{eq:irreducible}. The resulting polynomial has $L=29$.
For orientation, the retained counts
$i_2,\ldots,i_9$ are
$4,8,16,24,40,72,272,1200$.
Exact evaluation of \eqref{eq:lower-certificate} is negative at
\begin{equation}\label{eq:qminus}
 \qminus=\frac{3400034903}{10^9}=3.400034903.
\end{equation}
Thus $\mu>\qminus$. The notebook \texttt{mu-lower.ipynb}
generates all three raw tables, extracts the irreducible
coefficients, and checks this integer inequality.

\subsection{Upper bounds from finite windows}\label{sec:upper}

Fix $m\geq1$, and let $V_m$ be the set of valid words of
length $m$. Each valid word of length $m+1$ gives a directed
edge from its length-$m$ prefix to its length-$m$ suffix.
Let $A_m$ be the adjacency matrix with the \emph{column}
convention: $(A_m)_{v,u}$ is the number of edges from $u$ to $v$.
There are $S_{m+1}$ vertices and $S_{m+2}$ edges, obtained
directly from the rotation-word enumerator.
It follows that $\one^T A_m^t\one$ counts all paths of
$t$ edges, including their initial vertices.

\begin{theorem}\label{thm:window}
For $n\geq m+1$,
\begin{equation}\label{eq:upper-pointwise}
 S_n\leq \one^T A_m^{\,n-m-1}\one.
\end{equation}
Consequently $\mu\leq\rho(A_m)$, where $\rho$ denotes the
largest modulus of an eigenvalue.
\end{theorem}
\begin{proof}
By Lemma~\ref{lem:heredity}, sliding a window through a valid
word gives a path in this graph. The first vertex and each
subsequent appended symbol reconstruct the word, so the map
is injective. This proves \eqref{eq:upper-pointwise}.
For a finite nonnegative matrix, the sum of entries of $A^t$
has exponential rate $\rho(A)$: it is a matrix norm, up to
fixed factors, and Jordan normal form bounds its growth by
a polynomial times $\rho(A)^t$. Taking roots gives the claim.
\end{proof}

The graph enforces only bounded-window constraints. Its paths
may have collisions between wedges farther apart than the
window length, which is precisely why this construction gives
an upper, rather than an exact, count.

\subsubsection{Exact spectral certificates}
For any nonnegative matrix $A$ and positive vector $v$,
\begin{equation}\label{eq:collatz}
 Av\leq qv\quad\Longrightarrow\quad\rho(A)\leq q.
\end{equation}
Indeed, conjugating by $\operatorname{diag}(v)$ gives a
nonnegative matrix with every row sum at most $q$, whose
maximum-row-sum norm bounds all eigenvalues.
We use floating-point iteration only to propose $v$, then
round to a positive integer vector and check
$dAv\leq pv$ exactly for $q=p/d$.

The graph need not be strongly connected. Partition it into
\emph{strongly connected components}, the maximal sets in
which each vertex can reach every other. Ordering these
components makes $A_m$ block triangular, and its spectral
radius is the maximum of the diagonal-block spectral radii.
We retain the internal edges of every component containing a
cycle and discard singleton components without self-loops.
Call the resulting block diagonal matrix $A_m^{\rm cyc}$.
It has the same spectral radius as $A_m$.

On each retained component, iteration with $A+I$ and separate
normalization supplies a candidate vector. Adding $I$
removes periodic oscillation without changing the eigenvectors;
separate normalization prevents smaller components from
vanishing numerically. The final integer test
\eqref{eq:collatz}, not convergence of this iteration,
certifies the bound.

\begin{table}[t]
\caption{Finite-window certificates. States and edges refer to
$A_m^{\rm cyc}$, not the full graph. Endpoints are exact decimals.}
\label{tab:upper}
\centering
\begin{tabular}{r|r|r|c}
$m$ & States & Edges & Certified upper bound\\ \hline
9 & 153\,306 & 565\,050 & 3.685468366\\
11 & 2\,070\,446 & 7\,608\,583 & 3.674844805\\
13 & 27\,847\,372 & 102\,127\,233 & 3.667542939
\end{tabular}
\end{table}

Table~\ref{tab:upper} gives the finite computations in
\texttt{mu-upper.ipynb}. In particular, an integer-vector
certificate for $A_{13}^{\rm cyc}$ gives
\begin{equation}\label{eq:qplus}
 \mu\leq\qplus=\frac{3667542939}{10^9}=3.667542939.
\end{equation}
Together with \eqref{eq:qminus}, this proves the
computer-assisted interval
\[
 \boxed{3.400034903<\mu\leq3.667542939.}
\]

\subsection{Pointwise lower bounds}\label{sec:point-lower}

The renewal sequence \eqref{eq:renewal} already gives an exact
integer lower bound at each length. A pure exponential follows
from a finite induction check.

\begin{lemma}\label{lem:prefactor}
Suppose $\sum_{j=1}^L i_jq^{-j}\geq1$ and, for some $k_0$,
$c_k\geq aq^k$ for $k_0\leq k<k_0+L$.
Then $c_k\geq aq^k$ for every $k\geq k_0$.
\end{lemma}
\begin{proof}
For $k\geq k_0+L$ all predecessors $k-j$ are at least $k_0$.
Induction in \eqref{eq:renewal} gives
$c_k\geq aq^k\sum_j i_jq^{-j}\geq aq^k$.
\end{proof}

For the irreducible table above, take $L=29$, $k_0=2$,
$q=\qminus$, and $a=41/250$.
The root certificate supplies the required strict inequality
for the sum. The remaining 29 checks are the integer inequalities
\[
 250c_k(10^9)^k\geq41(3400034903)^k,\qquad 2\leq k\leq30.
\]
The notebook \texttt{pointwise-lower.ipynb} computes $c_k$,
verifies these checks, and provides a general prefactor
calculation for other cutoffs. It follows that
\begin{equation}\label{eq:closed-lower}
 S_n\geq\frac{41}{250}\qminus^{\,n-2}\qquad(n\geq4).
\end{equation}
Since $(41/250)\qminus^{-2}>0.014$ and
$\qminus>3.4000$, this implies
$S_n>0.014(3.4000)^n$ for $n\geq4$.
The exponential base is the same as in the lower bound
for $\mu$ because both conclusions use the same renewal
polynomial.

\subsection{Pointwise upper bounds}\label{sec:point-upper}

The full-graph expression \eqref{eq:upper-pointwise} is an
integer upper bound that can be evaluated by repeated sparse
matrix-vector multiplication. A simpler closed form is available
whenever a positive vector certifies the \emph{full} matrix.
If $A_mv\leq qv$, put
\[
 H=\frac{\sum_u v_u}{\min_u v_u}.
\]
Since $\one\leq v/\min_u v_u$, induction gives
\begin{equation}\label{eq:prefactor-upper}
 S_n\leq Hq^{\,n-m-1}\qquad(n\geq m+1).
\end{equation}

For $m=9$, the full graph has $172\,226$ vertices and
$633\,138$ edges. Its certificate in
\texttt{pointwise-upper.ipynb} has
\[
 \begin{gathered}
 q_9=\frac{3685468366}{10^9},\qquad \min v=1,\\
 \sum v=136209763558000711.
 \end{gathered}
\]
Thus
\begin{equation}\label{eq:closed-upper}
 S_n\leq136209763558000711\,q_9^{\,n-10}
 \quad(n\geq10).
\end{equation}
In the form $A_2q_9^n$, the prefactor satisfies
\[
 A_2=136209763558000711\,q_9^{-10}
 <294632981756.
\]
As $q_9<3.6855$, we obtain the rounded statement
$S_n<2.95\cdot10^{11}(3.6855)^n$ for $n\geq10$.
For short lengths the elementary upper bound $4^{n-1}$
is usually much smaller, and may always be used instead.

It is important not to insert the smaller base $\qplus$
into \eqref{eq:closed-upper} with this prefactor.
The componentwise $m=13$ certificate controls exponential
growth, but omits edges between components; such transitions
can contribute polynomial factors to path counts.
It is not a full-graph pointwise certificate.

\section{Computational reproducibility}\label{sec:reproduce}

All computations use integer lattice coordinates and the
collision rule of Section~\ref{sec:geometry}. The accompanying
code and notebooks are available in the repository
\begin{center}
\url{https://github.com/fedimser/math}
\end{center}
under \texttt{rubiks-snake/}. The four independent notebooks in
\texttt{asymptotic-analysis/} are:
\begin{itemize}
\item \texttt{mu-lower.ipynb}: enumerate slab blocks, apply
\eqref{eq:irreducible}, and certify \eqref{eq:lower-certificate};
\item \texttt{mu-upper.ipynb}: construct finite-window graphs
and check their componentwise spectral certificates;
\item \texttt{pointwise-lower.ipynb}: generate renewal counts
and certify the finite induction base;
\item \texttt{pointwise-upper.ipynb}: check a full-graph vector
and compute its exponential prefactor.
\end{itemize}
They depend only on the standard Python library, the shared
module \texttt{rubiks\_snake.py}, NumPy, Numba, and SciPy
(in addition to the notebook runtime). None reads unpublished
research files. Tests in \texttt{rubiks\_snake\_test.py}
include short-length comparisons of the block geometry
with rotation-word enumeration.

Each notebook exposes a bound function and evaluates successively
larger parameter choices. Increasing the slab cutoffs retains
more nonnegative irreducible coefficients; increasing the window
size enforces more local constraints. These changes strengthen
the mathematical approximations, although a finite-iteration
upper certificate may require more iterations or a finer
integer-vector scale to realize the improvement.
All substantive computations are timed. The archived largest
window run, $m=13$, used approximately 195 seconds and 6.5 GiB
of memory; these are historical resource measurements, not
requirements guaranteed on another machine.

The notebooks distinguish retained numerical outputs from newly
executed cells. Rerunning recomputes the relevant finite tables
and certificates, rather than treating stored output as input.
The reported decimals are rational bounds, not fitted values:
lower endpoints are accepted by an integer polynomial sign test,
and upper endpoints by an integer componentwise comparison.
Counter overflow is checked, and certificate products are
evaluated without fixed-width overflow. Floating-point iteration
is only a search procedure for a certificate; its error is not
part of the proof once the exact test passes.

\section{Conclusion}

The number of valid Rubik's Snake formulas has a well-defined
exponential growth constant. Exact finite computations, joined
to geometric concatenation and finite-window arguments, locate
it in $(3.400034903,3.667542939]$ and give explicit pointwise
exponential bounds. Complementary double occupancy and unique
slab factorization are essential: ignoring the former weakens
the lower family, while ignoring the latter can produce an
invalid lower bound.

Longer slab tables and larger window graphs can improve the
interval without changing the proofs. More efficient recognition
of forbidden local patterns may reduce the upper computation's
memory cost. A ratio limit, an algebraic correction factor,
or an equivalence $S_n\sim C\mu^n$ would require additional
arguments and is not established here.

\bibliographystyle{IEEEtran}
\bibliography{refs}
\end{document}